\section*{Capítulo \ref{ch:b2:heat-conduction} --- Conducción del calor}

\begin{solution}{exo:b2:heat-conduction:1}
$j = \lambda\Delta T/e = \qty{80}{W/m^2}$; \qty{8}{kW}; con lana, $R = 0.25 + 2.5 =
\qty{2.75}{m^2.K/W}$: \qty{7.3}{W/m^2}, \qty{730}{W}.
\end{solution}

\begin{solution}{exo:b2:heat-conduction:2}
$a$: \qty{1.2e-4}{m^2/s}, \qty{7e-7}{m^2/s}, \qty{1e-7}{m^2/s}. Para \qty{1}{cm}:
\qty{0.9}{s}, \qty{140}{s}, \qty{1000}{s}; para \qty{20}{cm}: \qty{6}{min}, \qty{15}{h},
\qty{5}{d}; corteza ($a \sim \qty{1e-6}{m^2/s}$): $10^{15}\,\unit{s}$, treinta
millones de años.
\end{solution}

\begin{solution}{exo:b2:heat-conduction:3}
Simple: $R = 0.169$, $U = \qty{5.9}{W/m^2/K}$, \qty{118}{W/m^2}; doble: $R =
0.63$, $U = \qty{1.6}{W/m^2/K}$, \qty{32}{W/m^2}. El hueco convecta y los vidrios
intercambian por radiación (cerca de la mitad de la transferencia real); el argón conduce
menos y convecta menos, y un tratamiento de baja emisividad suprime la
radiación.
\end{solution}

\begin{solution}{exo:b2:heat-conduction:4}
$L^2/a$ con el semiespesor: \qty{27}{min}; asado: $16\times$, \qty{7}{h};
$t \propto L^2 \propto m^{2/3}$: $\times 1.6$.
\end{solution}

\begin{solution}{exo:b2:heat-conduction:5}
(a) $(rT')'/r = -p/\lambda$ con $T'(0) = 0$. (b) $p = j^2/\gamma = \qty{1.7e6}{W/m^3}$;
$pR^2/4\lambda = \qty{1e-3}{K}$. (c) \qty{625}{K}. (d) Un hilo es isotermo; el centro de una
pastilla de combustible está cientos de kelvin por encima de su superficie --- la densidad
de potencia de un reactor la limita la conducción en el combustible.
\end{solution}

\begin{solution}{exo:b2:heat-conduction:6}
(a) \cref{prop:b2:heat-conduction:wave}. (b) \qty{12}{cm}; \qty{2.2}{m}. (c)
$\delta\ln 10 = \qty{5.2}{m}$; retraso $2.3\,\unit{rad}$, $130$ días. (d) La
amplitud de \qty{20}{K} debe reducirse a la mitad: $x = \delta\ln 2 = \qty{1.5}{m}$.
\end{solution}

\begin{solution}{exo:b2:heat-conduction:7}
(a) $\mathrm{Bi} = hR/3\lambda \approx 10^{-4}$; $\tau = \rho cR/3h = \qty{1100}{s}$.
(b) $\mathrm{Bi} = 1.2$: ningún límite; gobierna el tiempo interno $R^2/a \approx
\qty{6000}{s}$. (c) $R$ pequeño: $\tau$ minúsculo y \omterm{def:b2:heat-conduction:biot}{número de Biot} pequeño.
(d) Cae la resistencia de película $1/h$, que es la dominante.
\end{solution}

\begin{solution}{exo:b2:heat-conduction:8}
(a) \cref{prop:b2:heat-conduction:fin}. (b) $m = \sqrt{2h/\lambda t} =
\qty{13}{m^{-1}}$, $mL = 0.39$, $\eta = 0.95$. (c) Área $\times 2L/t = 60$, flujo
$\times 57$. (d) Más allá de $mL \approx 2$ la longitud sobrante está a la temperatura
del fluido; una aleta más fina tiene un $m$ mayor y un rendimiento menor.
\end{solution}

\begin{solution}{exo:b2:heat-conduction:9}
(a) \cref{prop:b2:heat-conduction:resistance}. (b) Lana $\ln 2/2\pi\lambda =
\qty{2.76}{K.m/W}$, película exterior $1/2\pi r_2h = 0.16$: \qty{45}{W/m}; desnuda,
$2\pi r_1h\Delta T = \qty{410}{W/m}$. (c) $\dd/\dd r[\ln(r/r_1)/2\pi\lambda + 1/2\pi hr]
= 0$ en $r = \lambda/h$. (d) $r_{\text{cr}} = \qty{2}{cm} > \qty{1}{mm}$: la funda
aumenta la pérdida --- enfría el cobre.
\end{solution}

\begin{solution}{exo:b2:heat-conduction:10}
(a) $8000(1/273 - 1/293) = \qty{2.0}{W/K}$. (b) $\sigma_s = \lambda T'^2/T^2$ con
$T' = \qty{100}{K/m}$; $\int\sigma_s\,\dd x = \lambda T'(1/T_{\text{fr}} - 1/T_{\text{cal}})
= j(1/T_{\text{fr}} - 1/T_{\text{cal}})$; por \qty{100}{m^2}: lo mismo. (c)
\qty{0.18}{W/K}. (d) $\Phi(1 - T_{\text{fr}}/T_{\text{cal}}) = \qty{550}{W}$; nada
--- el calor está ahora a $T_{\text{fr}}$ ($T_{\text{fr}}\dot S_{\text{c}} = \qty{550}{W}$).
\end{solution}

\begin{solution}{exo:b2:heat-conduction:11}
(a) $j_Q(0,t) = -\lambda\partial_xT = \lambda\theta_0\sqrt2/\delta\,\cos(\omega t +
\pi/4) = \sqrt{\lambda\rho c\omega}\,\theta_0\cos(\omega t + \pi/4)$. (b) El flujo de cada
lado en el contacto vale $b_i|T_i - T_{\text{cont}}|/\sqrt{\pi t}$; igualando:
$T_{\text{cont}} = (b_1T_1 + b_2T_2)/(b_1 + b_2)$. (c) Cobre \qty{20.3}{\celsius},
madera \qty{29}{\celsius}. (d) Baldosa $b \approx 1500$, alfombra $\approx 100$; la
fórmula vale mientras ambos cuerpos parecen semiinfinitos; una vez que los frentes
de calor alcanzan el riego sanguíneo y la cara lejana del objeto, el estado
estacionario depende del calor realmente suministrado.
\end{solution}

\begin{solution}{exo:b2:heat-conduction:12}
(a) Sustitúyase; $\theta(\pm L) = 0$ da $\cos k_nL = 0$. (b) $a_n =
4\theta_0(-1)^n/(2n+1)\pi$; $\theta = \sum a_n\cos(k_nx)\eu^{-ak_n^2t}$. (c) $n = 0$,
$\tau = 4L^2/\pi^2a$. (d) $\tau = \qty{650}{s}$; centro $(4/\pi)\eu^{-t/\tau} = 0.01$:
$t = 4.8\tau \approx \qty{52}{min}$.
\end{solution}

\begin{solution}{pb:b2:heat-conduction:1}
\textbf{1.} \qty{0.25}{m^2.K/W}; \qty{80}{W/m^2}.

\textbf{2.} $R = 0.415$, $U = \qty{2.4}{W/m^2/K}$; aislado, $R = 2.9$, $U =
\qty{0.34}{W/m^2/K}$.

\textbf{3.} \qty{5.9}{W/m^2/K}; \qty{1.6}{W/m^2/K}.

\textbf{4.} $4800 + 2360 = \qty{7.2}{kW}$; $680 + 640 = \qty{1.3}{kW}$.

\textbf{5.} $UA = \qty{358}{W/K}$: $358 \times 2500 \times 86400 = \qty{7.7e10}{J} =
\qty{21500}{kWh}$; después, \qty{66}{W/K}: \qty{4000}{kWh}.

\textbf{6.} $20 - 48/8 = \qty{14}{\celsius}$; $20 - 6.8/8 = \qty{19.2}{\celsius}$;
un muro frío radia menos hacia el cuerpo y puede alcanzar el punto de rocío.

\textbf{7.} $a = \qty{5.3e-7}{m^2/s}$; $e^2/a = \qty{7.5e4}{s} \approx \qty{21}{h}$: el
muro promedia el día (la $\delta = \qty{12}{cm}$ diaria es menor que
su espesor).

\textbf{8.} Búsquese $\theta_0\eu^{\iu(\omega t - \underline kx)}$: $\underline k = (1 - \iu)/\delta$.

\textbf{9.} \qty{12}{cm}; \qty{2.2}{m}.

\textbf{10.} $4.6\delta = \qty{54}{cm}$; \qty{3.7}{K}; \qty{1.1}{K}.

\textbf{11.} Un radián, $58$ días: mediados de marzo (siendo el mínimo de la
superficie a mediados de enero).

\textbf{12.} Amplitud $\lambda\theta_0\sqrt2/\delta = \qty{6.3}{W/m^2}$, adelantada
$\pi/4$ ($45$ días); cien veces el flujo geotérmico.

\textbf{13.} $\lambda \times \qty{0.03}{K/m} \approx \qty{0.03}{W/m^2}$, el orden correcto; son
\qty{0.1}{K} constantes en la altura de la bodega, sepultados en la
oscilación estacional.

\textbf{14.} \qty{1e6}{W/m^2}; $R = \qty{0.033}{K/W}$.

\textbf{15.} \qty{0.01}{K/W}; con aire: \qty{1.9}{K/W}, \qty{190}{K}.

\textbf{16.} $m = \qty{20}{m^{-1}}$, $mL = 0.61$, $\eta = 0.89$; $\mathrm{Bi} = ht/2\lambda
= 10^{-4}$.

\textbf{17.} $A = \qty{0.108}{m^2}$; $R_{\text{aletas}} = 1/\eta hA = \qty{0.21}{K/W}$; base
\qty{0.006}{K/W}; total \qty{0.26}{K/W}: \qty{26}{K}, unión a \qty{51}{\celsius}.

\textbf{18.} $1/hS = \qty{5.6}{K/W}$: \qty{560}{K} --- imposible.

\textbf{19.} $R_{\text{aletas}} \approx \qty{0.94}{K/W}$, total $\approx \qty{1}{K/W}$:
\qty{125}{\celsius} --- se ralentiza y luego se apaga.

\textbf{20.} $V = \qty{7.2e-5}{m^3}$, \qty{0.19}{kg}, $C = \qty{175}{J/K}$; $\tau = RC
\approx \qty{46}{s}$; pastilla: $C = \qty{0.07}{J/K}$, $\tau \approx \qty{3}{ms}$.

\textbf{21.} El calor debe repartirse de \qty{1}{cm^2} a \qty{36}{cm^2}: el cobre
reduce a la mitad esa resistencia de reparto; un caloducto mueve el calor por
evaporación y condensación con una conductividad efectiva cien
veces la del cobre.

\textbf{22.} $7200(1/273 - 1/293) = \qty{1.8}{W/K}$; \qty{0.33}{W/K}.

\textbf{23.} $100(1/298 - 1/324) = \qty{0.027}{W/K}$; $100/324 = \qty{0.31}{W/K}$ ---
diez veces más: convertir trabajo en calor es la gran irreversibilidad.

\textbf{24.} $\int\lambda T'^2/T^2\,\dd x = \lambda T'\,[-1/T] = j(1/T_{\text{fr}} -
1/T_{\text{cal}})$.

\textbf{25.} Muro: resistencias en serie, manda el peor conductor.
Aleta: $m = \sqrt{hp/\lambda S}$, $mL \approx 1$. Bodega: $\delta = \sqrt{2a/\omega}$,
amortiguamiento y retraso a la vez. Chip: una cadena de resistencias de la pastilla al
aire, con el área de las aletas haciendo el trabajo. En todas partes: el calor fluye cuesta abajo y
crea entropía al hacerlo.
\end{solution}
